\documentclass[12pt]{article}
\usepackage{amsmath}
\usepackage{amssymb}
\usepackage{amsthm}


\begin{document}
	\begin{enumerate}
		\item We need to check if $Q$ satisfies Kolmogorov's axioms
			\begin{itemize}
				\item $Q(A) = P(X \in A)$ so $Im(Q) \subseteq Im(P)$, therefore, $0 \leq Q(A) \leq 1 \quad \forall A \subseteq \mathbb{R}$.
				\item $Q(\emptyset) = P(X \in \emptyset) = P(X^{-1}(\emptyset)) = P(\emptyset) = 0$ and $Q(\Omega) = P(X \in \Omega) = 1$, since the value of $X$ has to correspond to an event which has to be a subset of the sample space. 
				\item (Countable additivity) \begin{align*} P\left(X \in \bigcup_{i \geq 1} A_i\right) &= P\left(X^{-1}\left(\bigcup_{i \geq 1}A_i\right)\right)\\
					&= P\left(\bigcup_{i \geq 1} X^{-1}(A_i)\right) \\
					&= \bigcup_{i \geq 1} P(X \in A_i)
					\end{align*}
					From analysis, function inverse commutes with unions (countable) and from $P$ satisfying the axioms already, we get the result. This is for any disjoint family of events $\{A_i\}_{i \in \mathbb{N}}$
			\end{itemize}
			\pagebreak
		\item Tiling are chosen based on a uniform random permutation of all possible coverings of the $2\times9$ board with $2\times1$ dominos. We must count how many possible coverings are there. By observation, there is only one way to cover the board such that all dominos are placed vertically. We have to count the total amount of ways of covering the board. 
			\begin{itemize}
				\item There is only one way of covering a $2\times 1$ board, call this $P_1 = 1$
				\item There are $2$ ways of covering a $2\times2$ board, call this $P_2 = 2$
				\item For any larger board $2\times n$, we have two ways of placing the topmost domino. If we place it vertically, then there are $P_{n-2}$ ways of covering the board, if we place it horizontally, then there are $P_{n-1}$ ways of covering the board.
			\end{itemize}
				 Thus $P_n = P_{n-1} + P_{n-2}$, this is the exact same thing as the Fibonacci recurrence with $P_0 = 0, P_1 = 1, P_2 = 2$, thus
				 \begin{align*}
					 P_1 &= 1 \\
					 P_2 &= 2 \\
					 P_3 &= 3 \\
					 P_4 &= 5 \\
					 P_5 &= 8 \\
					 P_6 &= 13 \\
					 P_7 &= 21 \\
					 P_8 &= 34 \\
					 P_9 &= 55
				 \end{align*}
				 Thus $P(\text{every domino is placed vertically}) = \frac{1}{55}$
				 \pagebreak
			 \item The probability that there is a vertical domino on the middle of the board, is $P_3 + P_4$, since there is a $2\times3$ board under and $2 \times 4$ board over the vertical domino tucked into a $2\times2$ square that has only one available covering, thus $P(\text{vertical domino on the middle of the board}) = P_3 + P_4 = \frac{3+5}{55} = \frac{8}{55}$
				 \pagebreak
			 \item 
				 The sample space here is $\Omega = S_n$ permutation group of $n$ numbers. We do this by fixing $\ell$ individuals, $i_1, i_2, \ldots, i_\ell$,  we first start by computing the probability that $i_1, i_2, \ldots, i_\ell$ gets the correct hat and nobody else does. From class, we know that the probability that no person , out of $n$ individuals, gets the correct that is : 
				 \[
					\sum_{k=1}^n \frac{ (-1)^k}{k!}
				 \]
				 We multiply this by $n!$, which is the size of the sample space (since it's a uniform random permutation), to get the number of arrangements where no person, out of $n$ individuals, gets the correct hat : 
				 \[
					n! \sum_{k=1}^n \frac{(-1)^k}{k!}
				 \]
				 Then we need to calculate number of arrangements such that $i_1, i_2, \ldots, i_\ell$ gets the correct hat and nobody else does. We fix $\sigma(i_j)$ for all $j \in [\ell]$, and consider the set of all other $n-\ell$ individuals. Thus the number of arrangements such that none of the other $n- \ell$ individuals gets the correct hat is:
				 \[
					(n-\ell)! \sum_{k=1}^{n - \ell} \frac{(-1)^k}{k!}
				 \]
				 Then we add over all $i_1 < i_2 < \ldots < i_\ell$, to get 
				 \begin{align*}
					 \sum_{1 \leq i_1 < i_2< \ldots < i_\ell \leq n}\left( (n-\ell)! \sum_{k=1}^{n - \ell} \frac{(-1)^k}{k!}\right) &= \binom{n}{\ell} (n-\ell)! \sum_{k=1}^{n - \ell} \frac{(-1)^k}{k!} \\
																	   &= \frac{n! (n-\ell)!}{\ell ! (n-\ell)!} \sum_{k=1}^{n - \ell} \frac{(-1)^k}{k!} \\
																	   &= \frac{n!}{\ell !}  \sum_{k=1}^{n - \ell} \frac{(-1)^k}{k!}
				 \end{align*}
				 To get the probability, we divide by the sample space $n!$ to get  
				 $$
					 P(\text{ $\ell$ people exactly get their hat and no one else does }) = \frac{ \sum_{k=1}^{n - \ell} \frac{(-1)^k}{k!}}{\ell !}
				 $$
				 Which in the limit as $n \to \infty$, as we saw in class, goes to 
				 $$
					\frac{e^{-1}}{l !}
				 $$
				 \pagebreak
			 \item 
				 Assume that there are uncountably many values of $\omega \in \Omega$ such that $\mathbb{P}(\{\omega\}) > 0$. 
				 Let $$E_n = \{ \omega \in \Omega \colon \mathbb{P}(\{\omega\}) > \frac{1}{n}\}$$
				 Then the $E_n$ must be finite for each $n$. Indeed, if the cardinality of $E_n$ is infinite, then we may construct a sequence of distinct points $\{\omega_k \colon k \in \mathbb{N}\} \subseteq E_n$, where the union of all points is a subset of $\Omega$, such that 
				 $$ \mathbb{P}\left( \bigcup_{k \in \mathbb{N}} \{\omega_k\} \right)  = \sum_{k \in \mathbb{N}} \mathbb{P}(\{\omega_l\}) > \sum_{k \in \mathbb{N}} \frac{1}{n} > 1$$ which raises a contradiction since $\mathbb{P}(E) \leq 1$ for any $E \subseteq \Omega$.
				 Then we claim that
				 $$ \bigcup_{n \geq 1} E_n = \{ \omega \in \Omega \colon \mathbb{P}(\omega) > 0\}$$ 
				 Indeed, if $\mathbb{P}(\omega) > 0 $ for some $\omega \in \Omega$, then $\mathbb{P}(\omega) > \frac{1}{N}$ for some $N \in \mathbb{N}$. It follows that $\omega \in \bigcup_{n \geq 1}E_n$. The converse is clear. But this raises a contradiction since
				 $$\bigcup_{n \geq 1} E_n $$ is a countable union of finite sets, which is at most countable. But $\{ \omega \in \Omega \colon \mathbb{P}(\omega) > 0\}$ is uncountable by our assumption. Therefore, we must conclude that the latter set is countable.
				 \pagebreak
			 \item 
				\begin{itemize}
				
					\item
	I'm assuming a hypothetical world where there are only two, binary, genders $ \{\text{boy}, \text{girl}\}$.  
						We already know that Miranda is giving birth to twins. Notice how we have events $Y = \text{identical twins}$ and $N = \text{not identical twins}$. These two latter events partition the sample space, since $\mathbb{P}(Y \cup N) = 1 \implies Y \cup N = \Omega$. 	We are looking for probability associated with the event $G = \text{ having two girls} \subseteq \mathcal{F}$, then using the law of total probability, noting that events $G$ and $I$, and $G$ and $N$ are independent. 
						$$\mathbb{P}(\text{having 2 girls}) = \mathbb{P}(G|Y)\mathbb{P}(Y) + \mathbb{P}(G|N)\mathbb{P}(N) = \frac{1}{3} = \frac{1}{2}\frac{1}{3} + \frac{1}{2} \frac{2}{3} = \frac{1}{2}$$
				 \item We are looking for 
					 $$\mathbb{P}(I|G) = \frac{\mathbb{P}(G|I) \mathbb{P}(I)}{\mathbb{P}(G)} =\frac{ \frac{1}{2} \frac{1}{3}}{\frac{1}{2}} = \frac{1}{3}$$
	 \end{itemize} 



	\end{enumerate}
	

	
\end{document}
